\scanpage{37}
\uline{NB} Si l'invariant $j$ est entier, alors on peut conjecturer que \deleted{la réduction de} $E_K$ a bonne réduction ssi il en est de $\unsure{\overline{\mathrm m}|=}$ $T_p(E_K)$, le grp $p$-divisible associé.
\margin{?}

c) L'invariant $j$ n'est pas entier.

Alors le \uline{modèle de \uline{Néron}} \inserted{$\unsure{\overline E|E}$} est tel que $E_k^0$ est une \inserted{forme de} $G_m$. Par suite, on conclut que dans ${}_{p^{\unsure{\nu|n}}}\unsure{\overline E|E}$ il y a un sous-groupe canonique \deleted{fini de} forme de $\mu_{p^{\unsure{\nu|n}}}$.
\margin{[qui n'est pas fini sur $A$ en gén]}
\figure{\begin{tikzpicture}\draw (0,-0.35) .. controls (0.15,0.2) and (0.15,0.2) .. (0.7,0.25)--(5,0.25) .. controls (5.3,0.25) and (5.4,0.3) .. (5.45,0.35);\end{tikzpicture}}

On trouve ainsi dans $T_p(E_K)$ un sous-groupe, qui est une forme \uuline{finiment} tordue de $T_p(G_m)$.
\margin{[donc un $T_p(G_m)$ défini au signe près]}
\figure{\begin{tikzpicture}\draw (0,-0.3) .. controls (0.1,0.2) and (0.2,0.2) .. (0.6,0.25) -- (4.5,0.4) .. controls (4.8,0.45) and (5,0.4) .. (5.2,0.25);\end{tikzpicture}}

Donc si $k$ est sép. clos, on trouve que $T_p(E_K)$ admet un sous-module isomorphe (de façon \uline{canonique} mod signe] à $T_p(G_m)\inserted{{}=Z_p(1)}$, quotient étant nécessairement isomorphe à $Z_p$ [car $\det T_p(E_K)\simeq Z_p(1)$].
\figure{\begin{tikzpicture}\draw[->] (0,-0.3) .. controls (0.15,0.2) and (0.8,0.1) .. (1.2,0.55);\end{tikzpicture}}

L'extension est ici déterminée par un élément canonique de $\widehat{K^*}$ ext. de $Z_p$ [\unsure{provenant|provient} de $K^*/A^*$] par $A^{**}$. Je conjecture
\margin{36}
